\chapter{\Swift}
\label{section:swift}


% \begin{figure}[htb]
 \begin{center}
  \includegraphics[width=0.2\textwidth]{61_pidt/ExCALIBUR.png}
 \end{center}


\noindent
\Swift\ is a SWIFT (SPH With Inter-dependent Fine-grained
Tasking) clone which combines the SWIFT SPH kernels
(tasks) with Peano's particle-in-dual-tree concept.
The original SWIFT codes is available from \url{https://swift.dur.ac.uk}.
\Peano's particle handling is a follow-up to the 2015 paper ``Two Particle-in-Grid Realisations on Spacetrees'' (\url{https://arxiv.org/abs/1508.02435}).
It is discussed in detail in Section \ref{section:pidt}.


The integration of the SWIFT kernels and \Peano\ has been made possible by the
EPSRC ExCALIBUR grants Massively Parallel Particle Hydrodynamics
for Engineering and Astrophysics (EP/V001523/1) and PAX-SPH.
\Swift\ is realised as Python layer on top of \Peano\ and brings \Peano's core, 
its particle toolsbox and SWIFT together.





\begin{framed}
\paragraph*{A first \Peano\ pidt run (native usage of particle interface)}
\begin{itemize}
  \item Configure \Peano\ with the flags \texttt{--enable-particles}:
  \begin{code}
libtoolize; aclocal; autoconf; autoheader; cp src/config.h.in .; automake --add-missing
./configure --enable-particles
  \end{code}
   These are the two mandatory extensions. You might want to add further options.
  \item Recompile your \Peano\ code. This should give you a couple of
  \texttt{ToolboxParticles} libraries in the \texttt{src/toolbox/particles}
  subdirectory.
  \item Ensure that \texttt{python} subdirectory is in your Python search path
  and
  \item Yet to be written
\end{itemize}
\end{framed}




\section{Set up a simulation}

Our interface for \Swift\ is almost complete realised in Python.
We configure the whole simulation in Python, but typically make the key 
steps of a simulation point to generic C++ routines which are shipped with
\Peano.
Alterantively, we can inject C++ code manually into the Python code (which is
great for prototyping).
Once the Python script is executed, it generates a plain C++ code with a
makefile which waves \Peano's core and the SPH kernels together.


We start with the creation of a \Swift\ project:
\begin{code}
project = swift2.Project( namespace=["benchmarks", "swift2", "planetorbit"], 
                          project_name="planetorbit", 
                          executable="orbit-exec" )
\end{code}


\marginpar{@todo Cristian}
\begin{itemize}
  \item global parameters
%   project.set_Peano4_installation(args.peanodir, build_mode)
% 
% project.set_load_balancing("toolbox::loadbalancing::RecursiveSubdivision", "new toolbox::loadbalancing::RecursiveSubdivision::DefaultConfiguration()")
  \item generate files
  \item compile
\end{itemize}

\section{Introduce a new particle}

\marginpar{@todo Critsian revise}
The generation of a particle is straightforward. We employ a tool called
DaStGen2 which really just is a set of Python classes which represent the 
attributes of a C++ class. We build up a set of Python objects and \Swift\
eventually dumps them into C++ code from which it is embedded into the 
actual simulator.   

\begin{code}
import dastgen2

particle = peano4.toolbox.particles.Particle("Particle")

v_attr = peano4.dastgen2.Peano4DoubleArray("v","Dimensions")
mass_attr = dastgen2.attributes.Double("mass")

particle.data.add_attribute(v_attr)
particle.data.add_attribute(mass_attr)


@todo something fancy in here that adds a new species
\end{code}

\begin{remark}
 Though conceptionally a plain C++ pretty printer, the DaStGen2 Python API
 can create some fancy C++ annotations. You can, for example, specify ranges
 in which integers live, or you can prescribe the precision you wanna have 
 for floating point values. If you have a specialised LLVM compiler, these
 annotations will be used to generate aggressively optimised code.
\end{remark}


\begin{definition}
 In \Swift's \emph{dual tree} storage scheme, particles are always stored to
 their nearest neighbour.
\end{definition}


\section{Input and output}

\Peano\ dumps a lot of information to the terminal by default. Consult
Section~\ref{section:logging:log-filter} for some remarks on how to filter the
output properly.


For the output of the data, the code can dump the particle data in various 
formats.
The easiest format to use is vtu data.

\marginpar{@Cristian Erklaerung}



% \subsection{todo}
% Python API muss
% \begin{itemize}
%   \item swift2.filter-file erzeugen und lesen
%   \item Load balancing offerieren
% \end{itemize}


\section{High-performance computing}
\label{section:particles:swift:hpc}

To use \Peano's high-performance computing features, you have to configure your
baseline code properly.
That is, you need to invoke \texttt{./configure} with an appropriate
\texttt{--with-multithreading=...} call to enable multithreading, 
you need to invoke it with \texttt{--with-mpi=...} for MPI support, and you need
to specify \texttt{--with-gpus=...} if you want to use accelerators.
Furthermore, you will have to enable \Peano's load balancing toolbox
(\texttt{--enable-loadbalancing}), unless you want to implement load balancing all on
your own.






\section*{Links and further reading}

\begin{itemize}
  \item The ``official'' pidt paper is
{\tiny \begin{verbatim}
@article{WEINZIERL201642,
title = "Two particle-in-grid realisations on spacetrees",
journal = "Parallel Computing",
volume = "52",
pages = "42 - 64",
year = "2016",
issn = "0167-8191",
doi = "https://doi.org/10.1016/j.parco.2015.12.007",
url = "http://www.sciencedirect.com/science/article/pii/S0167819115001635",
author = "T. Weinzierl and B. Verleye and P. Henri and D. Roose",
keywords = "Particle-in-cell, Spacetree, Particle sorting, AMR, Lagrangian–Eulerian methods, Communication-avoiding"
}  \end{verbatim}}
  If you use the software, it would be great if you could cite this one.
\end{itemize}


